% ============================================================
% Résumé — Français
% ============================================================

\chapter*{Résumé}
\addcontentsline{toc}{chapter}{Résumé}

La surveillance agricole dans les environnements sahariens présente des
caractéristiques distinctives qui la différencient des contextes agricoles classiques:
alimentation électrique intermittente, accès à Internet peu fiable, conditions
thermiques extrêmes, et difficulté pratique d'accéder aux parcelles éloignées entre
deux visites. Pour la culture de la luzerne en particulier, ces contraintes sont
accentuées par l'absence de systèmes de surveillance IoT conçus spécifiquement pour
cette culture dans le contexte agricole saharien, ainsi que par une limitation
récurrente identifiée dans les travaux examinés: la confusion entre l'horodatage de
téléversement et l'horodatage réel de la mesure, ce qui compromet la traçabilité de
toute analyse temporelle.

Cette thèse présente la conception et la mise en œuvre d'une plateforme de surveillance
IoT et d'analyse agronomique déterministe pour la culture de la luzerne dans la région
d'El Oued, en Algérie. Le système repose sur trois nœuds embarqués communicant par
liaison sans fil LoRa à 433~MHz selon une trame de dix minutes. Le nœud passerelle
principal coordonne la réception des paquets, maintient une horloge temps réel DS3231
de référence et enregistre toutes les mesures localement sur une carte SD,
indépendamment de la disponibilité du réseau. Le nœud Node2 mesure l'humidité, la
température et la conductivité électrique du sol au niveau du second pivot d'irrigation
via RS485/Modbus~; le nœud Node3 acquiert la température de l'air, l'humidité relative
et la pression atmosphérique. Le transfert des données suit un modèle axé sur le
stockage local: les enregistrements SD s'accumulent localement et sont transmis par
lot vers un serveur distant lorsque la connexion Internet est ponctuellement disponible.

Huit modes de fonctionnement du micrologiciel répondent aux contraintes spécifiques du
déploiement: un mécanisme de récupération corrige la dérive temporelle à l'origine
d'échecs de communication persistants, et des stratégies d'économie d'énergie
prolongent l'autonomie des batteries par sommeil profond et contrôle de l'alimentation
des capteurs par relais. Le serveur, développé avec Node.js et PostgreSQL, applique
une politique stricte qui sépare l'horodatage issu de l'horloge temps réel
(\texttt{measured\_at}) de l'horodatage de réception du téléversement. Un tableau de
bord React présente les tendances des capteurs, les événements agronomiques et les
résultats d'un pipeline déterministe en dix étapes de contrôle qualité et de scoring
de fiabilité. L'interprétation assistée par IA (Gemini) est encadrée par six règles
architecturales qui interdisent au modèle de calculer des métriques primaires ou de
modifier les données enregistrées.

Une fenêtre de validation pilote couvrant la période du 6 au 11 mai 2026 a produit
2433 relevés de capteurs répartis équitablement entre les trois nœuds, six sessions
de téléversement menées à terme, et 135 lignes issues du pipeline de traitement. Ces
résultats démontrent la faisabilité fonctionnelle de l'architecture de bout en bout
dans le cadre d'un déploiement pilote~; la validation sur une saison complète et
l'étalonnage des seuils de fiabilité constituent des priorités pour les travaux futurs.

\vspace{1em}
\noindent\textbf{Mots-clés:} Internet des Objets, agriculture intelligente, luzerne,
LoRa, ESP32, priorité locale, analyse déterministe, agriculture saharienne, qualité
des données, horloge temps réel, traçabilité temporelle.
